%%=============================================================================
%% Conclusie
%%=============================================================================

\chapter{Conclusie}%
\label{ch:conclusie}

% TODO: Trek een duidelijke conclusie, in de vorm van een antwoord op de
% onderzoeksvra(a)g(en). Wat was jouw bijdrage aan het onderzoeksdomein en
% hoe biedt dit meerwaarde aan het vakgebied/doelgroep? 
% Reflecteer kritisch over het resultaat. In Engelse teksten wordt deze sectie
% ``Discussion'' genoemd. Had je deze uitkomst verwacht? Zijn er zaken die nog
% niet duidelijk zijn?
% Heeft het onderzoek geleid tot nieuwe vragen die uitnodigen tot verder 
%onderzoek?

Deze bachelorproef bestudeerde de ontwikkeling en toepassing van een machine learning-methodiek die gebaseerd is op sensordata, om afwijkend gedrag en mogelijke storingen in bouwmachines te identificeren met als doel voorspellend onderhoud te ondersteunen. Aan de hand van de uitgevoerde experimenten kan worden vastgesteld dat een dergelijke aanpak technisch uitvoerbaar is, maar dat de betrouwbaarheid ervan sterk afhankelijk is van de kwaliteit van de beschikbare sensorgegevens, de aanwezigheid van historische defectinformatie en de gekozen machine learning-model.

De ontwikkelde methodiek omvat diverse opeenvolgende stappen. Eerst werden de real-world telematicadata van een operationele tractor verzameld en voorbereid. De verschillende databestanden werden bekeken, gestandaardiseerd en samengevoegd, waarna de belangrijke technische parameters werden gekozen. Daarna werd de privédataset bestudeerd om gebruikelijke operationele patronen en mogelijke anomalieën in de sensordata te achterhalen. Doordat de tractorgegevens geen historische defectlabels bevatten, werd er extra gebruikgemaakt van de Metro\-PT-3-dataset. Deze dataset zorgde ervoor dat de diverse modellen beoordeeld worden op basis van daadwerkelijk vastgelegde defectmomenten. Uiteindelijk wer\-d dezelfde methodiek gebruikt voor de private tractordataset om te bepalen of de modellen ook nuttig zijn voor de werkelijke sensordata van een bouwmachine.

In deze methodiek zijn vijf machine learning-modellen bestudeerd: Random Forest-model, Isoaltion Forest-model, Autoencoder, LSTM Autoencoder en One-Class SVM-model. De uitkomsten laten zien dat de kezuze van het model een cruciale impact heeft op de bruikbaarheid van de vastgelegde afwijkingen. Random Forest-model bereikte bijvoorbeeld een uiterst hoge accuracy, maar slaagde er nauwelijks in om de vastgelegde fouten te identificeren. Dit laat zien dat accuracy alleen geen geschikte indicator is voor een zeer onevenwichtige dataset. Isolation Forest-model had ook een beperkte detectie van de bekende tekortkomingen. Binnen de bestudeerde modellen liet de Autoencoder verbeterde resultaten zien bij het opsporen van de defectklasse met een recall van 0.44 en een F1-score van 0.24 tijdens de aangepaste MetroPT-3-testperiode. De LSTM Autoencoder verbeterde niet in vergelijking met de gewone Autoencoder. De One-Class SVM-model boekte wel een hoge recall, maar beschouwde tegelijkertijd een aanzienlijk aantal normale observaties als een anomalie.

Deze uitkomsten geven op een genuanceerde wijze een antwoord op de onderzoeksvraag. De machine learning-methodiek die is onderzocht, kan dienen om afwijkingen in industriële sensordata van bouwmachines te herkennen. Vooral unsupervised en reconstructiegebaseerde technieken zoals de Autoencoder kunnen mogelijkheden bieden voor situaties waarin er geen defectlabels beschikbaar zijn. De aanpak kan fungeren als fundament voor een systeem dat onderhoudsverantwoordelijken waarschuwt voor sensorgedrag dat niet overeenkomt met het aangeleerde gebruikelijke gedrag. Aan de hand van de huidige uitkomsten kan nog niet worden vastgesteld dat een vastgelegde afwijking automatisch gelijkstaat aan een echte technische storing.

Een cruciale rol van deze bachelorproef is de creëring en beoordeling van een complete data-gedreven methodiek, variërend van het verzamelen en voorbereiden van real-world telematicagegevens tot het trainen, beoordelen en vergelijken van diverse machine learning-modellen. De mix van de private tractordataset en de MetroPT-3-dataset zorgde ervoor dat zowel de praktische bruikbaarheid van de methodiek als de prestaties van de modellen bij bekende defecten te evalueren. Een dergelijke methodiek kan in de toekomst kleine bouwbedrijven helpen bij het monitoren van machinegedrag en het vroegtijdig opmerken van mogelijke afwijkingen zonder dat een defect eerst moet optreden voordat onderhoud plaatsvindt. Dit heeft niet alleen de technische effectiviteit van specifieke modellen onderzocht, maar ook de haalbaarheid van een complete anomaliedetectie-methode in een realistische industriële situatie.

De uitkomsten  voldeden gedeeltelijk aan de verwachtingen. Vooraf was het duidelijk dat modellen die speciaal zijn ontwikkeld voor het opsporen van afwijkingen beter zouden passen bij de privédataset, aangezien er geen gelabelde defecten in waren. De experimenten gaven echter aan dat een ingewikkelder model niet automatisch resulteert in betere uitkomsten. De LSTM Autoencoder presteerde niet beter dan de gewone Autoencoder, ondanks het gebruik van tijdelijke informatie. Daarnaast werd aangetoond dat een zeer hoge recall, zoals bij One-Class SVM-model, niet genoeg is als dit samengaat met een aanzienlijk aantal foutieve alarmen. Dit onderstreept hoe cruciaal het is om samen de precision, recall, F1-score en de praktische waarde van de vastgestelde afwijkingen te evalueren.

De voornaamste beperking van de methodiek blijft het gebrek aan bevestigde defectlabels in de private tractordataset. Hierdoor kunnen de vastgelegde afwijkingen niet direct worden verbonden met echte defecten of onderhoudswerkzaamheden. De methodiek dient momenteel te worden gezien als een ondersteunend detectiesysteem in plaats van als een volledig zelfstandig voorspellend onderhoudssysteem. Een afwijking toont aan dat het sensorgedrag niet overeenkomt met het gedrag dat het model als normaal heeft aangenomen, wat betekent dat er aanvullend technisch onderzoek vereist is.

Het onderzoek resulteert hierdoor in diverse opties voor toekomstig onderzoek. Een langere historische meetperiode, verschillende bouwmachines en een gestructureerde registratie van onderhouds- en storingsmomenten zouden het mogelijk maken om de methodiek grondiger te verifiëren. Daarnaast kan worden bestudeerd welke modelconfiguraties en anomaliedrempels het beste balans creëren tussen het op tijd opsporen van storingen en het verminderen van onjuiste alarmen.

Deze bachelorproef toont dat het mogelijk is om een data-gedreven machine lear\-ning-methodiek toe te passen voor het opsporen van afwijkend gedrag in sensordata van bouwmachines. De methodiek kan voornamelijk worden gebruikt als hulpmiddel voor voorspellend onderhoud, maar het is essentieel om deze verder te valideren met langere en gelabelde praktijkdata voordat ze betrouwbaar kan worden toegepast voor echte voorspelling van storingen. Het onderzoek biedt een eerste praktische fundament voor het toepassen van machine learning gebaseerde detectie van afwijkingen in het onderhoud van bouwmachines.









